% Editable editorial notes, separate from the transcribed book.
\EditorialNotes

Felix Hausdorffs Werk \emph{Sant' Ilario. Gedanken aus der Landschaft Zarathustras}
erschien 1897 in Leipzig bei C. G. Naumann unter dem Pseudonym Paul Mongré.
Eine Einordnung dieses sehr eigentümlichen Textes sowie eine Biographie des Autors
finden sich in Band VII der \emph{Gesammelten Werke} von Felix Hausdorff, herausgegeben im Springer-Verlag durch Werner Stegmaier.
Dort finden sich auch nützliche Zeilenkommentare zum Text.

Felix Hausdorff nahm sich im Jahr 1942 das Leben.
Nach Ablauf der Schutzfrist Ende 2012 ist sein Originaltext seit dem 1. Januar 2013 in Deutschland gemeinfrei.
Da es bis auf den Facsimile-Nachdruck in den \emph{Gesammelten Werken} keinen weiteren Nachdruck gibt, habe ich mir die Mühe gemacht,
eine digitale Neusetzung auf Grundlage der Leipziger Ausgabe zu erstellen.

Für diese Neusetzung wurde die Originalausgabe gescannt und der Text mit Hilfe verschiedener OCR-Werkzeuge (hauptsächlich auf Grundlage von LLMs) transkribiert.
Die Transkription wurde anschließend durch eine Kombination automatischer und händischer Korrektur bearbeitet.

\EditorialHeading{Textgestalt}
Ich habe mich bemüht, weitgehend das Schriftbild des Originals zu erhalten.
Insbesondere wurde die ursprüngliche Rechtschreibung und die Verwendung von Sperrsatz,
um Wörter hervorzuheben, beibehalten.

\EditorialHeading{Fremdsprachige Wörter und Wendungen}
Die Kursivierung fremdsprachiger Ausdrücke ist eine editorische Zutat.
In vielen Fällen erleichtert sie es dem Auge des Lesers, sich im Text zurechtzufinden.
Allerdings ist die Einschätzung eines Ausdrucks als ``Fremdwort'' in vielen Fällen nicht eindeutig.
Ich habe hier nach meinem eigenen Geschmack entschieden.
\par
Bei Substantiven gilt die im Text verwendete Großschreibung als
Zeichen ihrer Behandlung als deutsches Wort: ``Façon'' etwa bleibt deshalb
aufrecht. Hingegen wird \Foreign{point d’honneur}
kursiv gesetzt, auch wenn die Wendung zum vertrauten Sprachgebrauch
der Zeit gehört.
\par
\ifEditorialItalics
In dieser Leseausgabe ist die zusätzliche Kursivierung verwendet.
\else
In dieser Leseausgabe ist die zusätzliche Kursivierung ausgeschaltet.
\fi
\par

\EditorialHeading{Einzelne Eingriffe}
Im Original wird der Aphorismus zwischen 223 und 225 irrtümlich als „124“ abgedruckt,
dieser Satzfehler ist hier korrigiert worden.
\par

\EditorialHeading{Zur digitalen Ausgabe}
Neusetzung und editorische Ergänzungen: Arne Hofmann.
Erste digitale Ausgabe, 3. November 2026.
Zeilen- und Seitenumbrüche sind gegenüber der Originalausgabe neu gesetzt.
\par
Diese Leseausgabe ist keine historisch-kritische Edition.
Verbliebene Übertragungsfehler sind möglich.
\par
Neue editorische Beiträge, soweit urheberrechtlich geschützt:
\copyright\ 2026 Arne Hofmann, \mbox{CC BY-NC 4.0}.
\par
\url{https://creativecommons.org/licenses/by-nc/4.0/}
\par
Die Lizenz erfasst weder den gemeinfreien Originaltext noch die eingebetteten
Schriften; für diese gelten ihre eigenen Lizenzen.
